\section{A terminating exact-minimum procedure with explicit cycle input}
\label{sec:exact-algorithm}

A uniform upper certificate and an exact minimum are different outputs.
This section gives a finite procedure for the second output once the
arithmetic variety, target cycle, and bounded parameter spaces are
actually presented. Its termination uses the finite-field Tate theorem
accepted as an input of this task, specifically
\cite[Theorem 3.1 and \S3.2]{MilneReport2026}. It supplies no practical
running-time estimate.

\paragraph{Input convention.}
Let $X/\mathbb F_q$ be an explicitly presented abelian fourfold, let
$H$ be an explicit ample integral polarization, and let $C\geq1$ be an
integer known to satisfy $c_H(t)\leq C$. The target is given by actual
cycles with equations,
\[
             t=\frac{\cl(Z_0^+)-\cl(Z_0^-)}{n_0},
             \qquad n_0\in\mathbb Z_{>0}.
\]
The supports and their multiplicities may be defined over a specified
finite extension of $\mathbb F_q$. In addition, the complete
individual-cycle bounded-degree parameter space $\mathcal Q_C$ of B01
(the factor of its pair space $\mathcal H_C=\mathcal Q_C\times\mathcal Q_C$) is supplied by
explicit equations, with its universal families or effective
point-to-cycle maps, including every multiplicity type and the empty
cycle. Thus every effective surface cycle of $H$-degree at most $C$
occurs in this parameter space. Merely naming an abstract Betti or
$\ell$-adic vector is not this input model. Nor does the theorem silently
compute a Hilbert-polynomial bound that has not been supplied.

\begin{lemma}[R40a: exact computation of the intersection tests]
\label{lem:R40a}
For explicit surface cycles $Z,\beta$ on the smooth fourfold $X$, the
integer $\deg(Z\cdot\beta)$ is computable over a common finite field of
definition, including when the supports do not meet properly.
\end{lemma}
\begin{proof}
It is enough to treat integral supports and extend bilinearly over their
integer multiplicities. On a smooth variety the coherent sheaves
$\mathcal O_Z,\mathcal O_\beta$ are perfect. The leading term of the
Chern character of either sheaf is its fundamental codimension-two
cycle. The Hirzebruch--Riemann--Roch formula
\cite[Corollary 15.2.1, p.~288]{Fulton}, applied to their derived tensor
product, therefore gives
\[
 \deg(Z\cdot\beta)
 =\chi\bigl(X,\mathcal O_Z\mathbin{\otimes^{\mathbf L}}
                         \mathcal O_\beta\bigr).
\]
Indeed their Chern characters start in codimension two, so only the
product of the two leading terms can contribute in dimension four;
positive-degree Todd terms and higher Chern-character terms contribute
nothing. Finite locally free resolutions, or finite resolutions on a
projective affine cover followed by coherent-sheaf computations, compute
the Euler characteristic exactly. Equivalently one can compute the
finitely many Tor sheaves and their Hilbert polynomials. These are finite
polynomial-algebra operations. The degree is taken over a common field
of definition, so extending that field does not multiply the geometric
intersection number.
\end{proof}

\begin{theorem}[R40: an exact local minimum from finite intersection data]
\label{thm:R40}
Under the stated input convention and the cited Tate input, there is a
terminating algorithm that returns $c_H(t)$ and a representative
attaining it in one fixed prime-to-characteristic realization. Under
the supplied identification of numerical and homological equivalence
in each realization under consideration, the same algorithm returns
the simultaneous-realization minimum.
\end{theorem}
\begin{proof}
We first obtain a finite complete system of tests for equality of cycle
classes, and then search the bounded parameter space.

\emph{The geometric Tate dimension.}
Frobenius on $H^1$ of an abelian variety over a finite field is
semisimple: its rational algebra is contained in the center of the
semisimple endomorphism algebra. The argument appears in
\cite[\S2, pp.~138--139; Theorem 2(a), p.~140]{Tate1966}.
Its exterior powers are semisimple as well. The eigenvalues on
$H^4(X_{\overline{\mathbb F}_q},\mathbb Q_\ell(2))$ are the products
of four $H^1$ eigenvalues, divided by $q^2$.

All the relevant polynomials are computable from finite-field point
counts. For a deliberately inefficient finite method, the known Betti
numbers $b_i=\binom8i$ show that the numerator and denominator of the
zeta function each have degree $128$. Point counts through extension
degree $256$ determine its Taylor series through that degree and hence
the rational function by rational-function reconstruction. The Weil
weights separate its $H^1$ factor from the other factors. Taking the
fourth exterior power then computes the polynomial on $H^4(2)$.
Every point count in this prescription is itself a finite enumeration;
no complexity estimate is intended.

Determine which roots of the degree-$70$ polynomial on $H^4(2)$ are
roots of unity, retaining their multiplicities, and denote their total
multiplicity by $r$. This is a finite algebraic-number operation. For
example, a root of unity of order $m$ occurring here has
$\varphi(m)\leq70$; the elementary bound
$\varphi(m)^2\geq m/2$ reduces the test to $m\leq9800$.
Factorization and comparison with these cyclotomic polynomials suffice.
The resulting number $r$ is the dimension of the geometric Tate
subspace $V$, the union of the subspaces fixed by finite powers of
twisted Frobenius.

The Poincar\'e pairing is nondegenerate on $V$. After diagonalizing
Frobenius, an eigenvalue can pair only with its inverse. Roots of unity
have root-of-unity inverses, and no non-root-of-unity eigenspace can
pair with $V$. Nondegeneracy on the full middle cohomology therefore
restricts to nondegeneracy on $V$.

\emph{An algebraic basis and its intersection matrix.}
Enumerate effective integral surface cycles over all finite extensions
of $\mathbb F_q$, for example by enumerating homogeneous ideals and
their reduced irreducible supports. Compute their pairwise intersection
numbers by Lemma~\ref{lem:R40a}. Continue through finite prefixes until
the intersection matrix has rank $r$, and then choose $r$ cycles
$\beta_1,\ldots,\beta_r$ with nonsingular Gram matrix.
One must retain the whole finite prefix when testing rank: greedily
discarding cycles with self-intersection zero would be invalid.

This search terminates. Every enumerated cycle is defined over a finite
extension and hence has class in $V$. The Tate input says that these
classes span $V$, and its pairing is nondegenerate. Some finite prefix
therefore has rank $r$. The chosen cycles form a basis of $V$ over
$\mathbb Q_\ell$. For any cycle class $x\in V$, the rational vector
\[
                   v(x)=\bigl(\deg(x\cdot\beta_i)\bigr)_{i=1}^r
\]
is consequently a complete equality test: $v(x)=0$ if and only if
$x=0$ in the selected realization. Rational coefficients cause no
problem because the Gram determinant is a nonzero rational number.

\emph{The finite list of bounded cycle classes.}
Compute the finitely many geometric connected components of
$\mathcal Q_C$. Absolute decomposition and the incidence relations of
irreducible components are finite polynomial-algebra operations over
the algebraic closure of a finite field. Choose a closed point on each
component, by search over finite extensions if necessary, and use the
supplied universal families to obtain its effective cycle $Z_j$.
Record its $H$-degree $d_j$ and the exact intersection vector
$v_j=v([Z_j])$.

The cycle class is constant on each geometric connected component:
the ambient variety $X$ is fixed, and the relative class in the constant
cohomology local system is locally constant, as in B01. Thus this finite
list includes the class of every effective cycle of degree at most
$C$. It is harmless if several components give the same class. Include
$(d_j,v_j)=(0,0)$ for the empty cycle. If the parameter-space embedding
uses a multiple $kH$, its surface degree is divided by $k^2$ when
recording the $H$-degree.

\emph{The exact finite minimization.}
Compute
\[
 v(t)=\frac{v([Z_0^+])-v([Z_0^-])}{n_0}
\]
by exact rational arithmetic. For each integer $1\leq N\leq C$ and
each ordered pair of labels $i,j$, test
\[
                         v_i-v_j=Nv(t).
\]
Among the successful tests return the smallest
$\max(N,d_i,d_j)$, together with its two chosen cycles. There is at
least one successful test because $c_H(t)\leq C$. Every successful
test is an actual cohomological representation by the nondegenerate
basis test. Conversely, every representation of complexity at most
$C$ occurs in the list of component classes and multipliers. Hence
the returned number is exactly the minimum, not merely another upper
bound.

Finally, a zero intersection vector says that the difference is
numerically trivial. Under the supplied equality of numerical and
homological equivalence for every realization in use, this is the same
simultaneous equality. In particular the input target is an actual
rational cycle throughout; no oracle for rationality of an arbitrary
$\ell$-adic Tate vector has been introduced.
\end{proof}

\begin{corollary}[R41: a finite description of the entire local function]
\label{cor:R41}
Fix an arithmetic abelian fourfold with its extending polarization and
a compatible target class. Suppose that an actual generic rational
cycle representing the target, its exact generic minimum $C_0<\infty$,
and the integer $b_{C_0-1}$ of B04 have been computed. Suppose also that
the required bounded parameter spaces and universal families are
explicit at the remaining good places. Then the function
$v\mapsto c_v(t_v)$ on geometric good reductions is computable as the
constant value $C_0$ together with a finite table of exceptions.
\end{corollary}
\begin{proof}
If $C_0=1$, specialization gives $c_v(t_v)\leq1$ and the positive
multiplier convention gives equality, so no $b_0$ is needed.
For $C_0>1$, B03 and B04 give $c_v(t_v)=C_0$ outside their excluded
rational characteristics and the prime divisors of $b_{C_0-1}$.
There are only finitely many places of the fixed number field above
this finite set of rational primes. Retain those places at which the
given model has good reduction; bad-reduction places are outside the
domain of this function.

At each remaining good place, specialize the explicit rational cycle
to obtain the input target. Properness of the relevant cycle parameter
spaces, or the usual effective cycle specialization after a finite
extension, gives the upper bound $C_0$ from a generic representative
attaining that bound. Apply Theorem~\ref{thm:R40} with $C=C_0$ to
compute the exact value at this place. Residue-field extension does not
alter the geometric minimum. The finitely many resulting values,
together with the constant outside the listed places, give the asserted
finite description.
\end{proof}

This is a conditional computational consequence of explicit inputs,
including the generic minimum and the bounded parameter-space data in
R41. The abstract family bound R24 by itself supplies none of those
equations or target-cycle coordinates. Computing a generic minimum
from a bare Betti vector would require a separately specified exact
topological input model and comparison algorithm; no such additional
algorithm is asserted here.
